% ══════════════════════════════════════════════════════════════════════════════
\chapter{Conclusion and Future Work}
\label{ch:conclusion}
% ══════════════════════════════════════════════════════════════════════════════

\section{Research Questions Answered}

\textbf{Main research question.} Can an economical, portable SatNOGS-compatible ground station be designed, built, and field-validated as the ground segment for the AetherSpace CubeSat mission? \textbf{Yes, with one open item.} The station was designed, fabricated, and field-deployed across three sessions totalling approximately 210 tracked passes. The rotator meets its pointing requirements, the RF chain is end-to-end verified in ground-to-ground tests, SatNOGS network integration is in place via EasyComm III and SiDS, and the station deploys in approximately four minutes from a camera tripod. The single open item is orbital frame reception: no protocol-valid frames were decoded from orbit, with the root cause identified as a missing front-end LNA. The hardware change required to close the sensitivity gap is specified in \S\ref{sec:future-work}; once applied, the station is ready to receive AetherSpace telemetry when the satellite launches.

\begin{enumerate}
    \item Under €50 rotator / €100 station? Partially. v2 (\S\ref{sec:cost}) reaches $\sim$€40 for the combined controller/RF PCB, $\sim$€75 for the full rotator hardware, and $\sim$€100 for the station core (excluding antenna, coax, power supply, tripod). The strict ``full rotator under €50'' target is not met once motors and mechanics are included. The v1 setup used for field validation costs $\sim$€145 (two separate boards, 6-layer RF HAT).
    \item Pointing accuracy? Measured: $<$0.5\degree\ EL and 0\degree--3\degree\ AZ in the 20\degree--70\degree\ peak-elevation band (\S\ref{sec:tracking-accuracy}). Encoder resolution 0.00409\degree\ AZ, 0.00793\degree\ EL (\S\ref{sec:gearing}). Both axes stay inside the Siretta Oscar 44 3-dB beamwidths (54\degree\ H, 48\degree\ E). The large AZ excursions at $>$70\degree\ elevation are a motor-speed-limit event, not a control-loop accuracy limit.
    \item SatNOGS integration? Yes on the rotator side via standard EasyComm III / hamlib model 204, and on the telemetry side via SiDS (station 4712 registered). The CC1200 is incompatible with the SDR-based \code{satnogs-client}, so the station does not appear ``online'' on the observation-scheduling map; this is a deliberate choice, not a defect.
    \item DC vs stepper? DC motors were chosen (rationale in \S\ref{sec:motors}). Empirical results: zero motor current draw when stationary, $<$0.5\degree\ EL / 0\degree--3\degree\ AZ tracking in the link-closing elevation band, 35\degree/s peak slew at 95\,\% duty. Cost: 100~Hz PID with quadrature decode, running within the RP2040 timing budget (ISR latency $<$10~$\mu$s; inter-edge period 117~$\mu$s at 35\degree/s).
    \item CC1200 sensitivity? The assembled RF HAT achieves approximately $-$105~dBm in bench conditions and $-$96~dBm outdoors (\S\ref{sec:packet-reception}). The dominant limiting factor is board-level interference: LMR51420 buck-converter harmonics, RP2040 digital noise, and Pi GPIO emissions lift the effective noise floor approximately 30~dB above the thermal\,+\,NF baseline, degrading sensitivity from the datasheet $-$120~dBm to approximately $-$105~dBm and leaving the AetherSpace link at $-$20.7~dB measured margin at 30\degree\ elevation. A mast-mounted LNA is the single hardware change required to recover the lost margin (\S\ref{sec:future-work}).
    \item CC1200 vs SDR? For AetherSpace specifically, the CC1200 approach is the better fit: both ends use the same transceiver, so a single SmartRF register set loaded at both ends eliminates modulation-parameter uncertainty. The trade-off is reduced flexibility for other satellites and incompatibility with the standard \code{satnogs-client} pipeline. The sensitivity gap documented in Q5 is not inherent to the CC1200 approach; it is a board-interference problem that an LNA would resolve equally for any narrow-band receiver on this hardware.
\end{enumerate}

\section{Limitations}

\begin{itemize}
    \item No protocol-valid frames decoded from orbit: documented in \S\ref{sec:packet-reception} with per-observation alternative hypotheses; root cause identified in \S\ref{sec:link-budget} and \S\ref{sec:field-tests} as the absence of an external LNA combined with a board-noise floor approximately 30~dB above the expected kTB+NF limit, degrading effective CC1200 sensitivity from the datasheet $-$120~dBm to approximately $-$105~dBm and leaving the AetherSpace link at approximately $-$20~dB measured margin at 30\degree\ elevation (best case; typical outdoor conditions are worse).
    \item No calibrated-source RF verification against a signal generator. Sensitivity has been validated in relative terms (ground-to-ground CC1200-to-CC1200 decode, \S\ref{sec:packet-reception}) but not against a signal generator injecting a known power level into the SMA. A sig-gen sweep would separate ``receive chain works but insensitive'' from ``receive chain has an undiagnosed defect'' independently of link geometry.
    \item VHF path not validated. The 144/145~MHz front-end is routed on the RF HAT but not populated. It would need a VNA-based matching-network tune before uplink use.
    \item Zenith AZ rate saturation. The motor's 35\degree/s ceiling is exceeded briefly during high-elevation passes (near-zenith passes can demand $>$10\degree/s of \textit{azimuth} rate for a few seconds, even though the total angular rate never exceeds $\sim$1.1\degree/s). The AZ axis saturates during this window and recovers as the satellite descends below $\sim$70\degree\ elevation.
    \item First-rev PCB issues: bodge wires (AZ IN1 rerouted GP15$\to$GP16), replaced TB6642FG, removed encoder RC caps. All addressed in the v2 MOTOR\_RF\_HAT design.
    \item Design-complete v2 not yet built. The cost claim relies on v2 BOM pricing; v2 has not been fabricated in time for this thesis. The v2 design incorporates every lesson from v1 bring-up and is ready to order.
\end{itemize}

\section{Future Work}
\label{sec:future-work}

\begin{enumerate}
    \item Build the v2 MOTOR\_RF\_HAT: fabricate and assemble the design-complete combined board. Validates the $\sim$€100 station-core cost projection end-to-end and removes the v1 bodges.
    \item External mast-mounted LNA: the single most impactful reception change. Concrete recommendation:
    \begin{itemize}
        \item LNA: Qorvo QPL9547 (NF $\sim$0.3~dB, $\sim$20~dB gain at 433~MHz, +5~V) as a modern replacement for the discontinued SPF5189Z. Mini-Circuits PGA-103+ (NF $\sim$0.6~dB, $\sim$20~dB gain at 433~MHz) is an equivalent in-stock part. There is no TI companion LNA for this band; the CC1190 range extender covers 850--950~MHz only.
        \item SAW pre-filter (optional but recommended): 433.92~MHz SAW bandpass before the LNA, to reject out-of-band ISM / PMR / ATV interferers that otherwise compress the LNA and lift the AGC.
        \item Architecture: mast-mount the LNA, fed over the antenna SMA coax via a bias-T on the RF HAT (a 470~nH--1~$\mu$H RF choke plus DC-block capacitor; the choke reactance must be well above $Z_0 = 50~\Omega$ at 433~MHz to keep insertion loss negligible: 470~nH gives $X_L \approx 1.3~\text{k}\Omega$, yielding ${\approx}$0.2~dB loss; 120~nH gives only 326~$\Omega$ and ${\approx}$0.6~dB, which is significant ahead of a low-noise amplifier). Placing the LNA at the antenna feed rather than after the coax makes the LNA's NF set the whole-chain NF (Friis cascade) and pre-amplifies before the 0.5--1~dB RG-58 loss and the 15--20~dB board-noise elevation measured on v1.
        \item Antenna (optional secondary upgrade): a 10--12~dBi Yagi in place of the Oscar~44 adds a few dB, but the dominant gap is the LNA and the antenna upgrade is not required to close the link.
    \end{itemize}
    Combined impact (LNA alone, existing Yagi): system NF $\sim$7~dB $\to$ $<$1~dB, AGC lifts off the board-noise floor. Effective sensitivity improves by 15--20~dB, moving the AetherSpace link at 30\degree\ elevation from $-$20.7~dB measured to roughly 0\ldots$+$5~dB. v2 PCB change: add a bias-T network on the UHF SMA port; the LNA itself stays off-board.
    \item AetherSpace integration: once AetherSpace's final packet format is published, load the matching SmartRF profile. A CC1200-to-CC1200 link removes a large part of the modulation-parameter uncertainty if both ends use the same register profile; the open question is sensitivity, addressed by (2).
    \item RTL-SDR hybrid: add an RTL-SDR for wideband IQ capture and visual waterfall. Useful as an independent verification path for CC1200 reception claims. Software architecture already supports it via the abstract \code{RfBackend} interface.
    \item VHF uplink: populate the 144/145~MHz matching network, VNA-tune, and enable uplink commanding once AetherSpace defines an uplink protocol.
    \item Battery field test: validate the 20--28~h runtime projection with current measurements at each operating point on a representative 6S pack.
\end{enumerate}

\section{Summary}

This thesis designs, builds, and partially validates a portable SatNOGS-compatible ground station. The v1 hardware (two separate boards, used in the field sessions) tracks satellites with sub-0.5\degree\ elevation error and 0\degree--3\degree\ azimuth error in the 20\degree--70\degree\ peak-elevation band, integrates with the SatNOGS rotator ecosystem through standard EasyComm III, and deploys from unpacked to tracking in about four minutes using only a laptop (a phone can also supply GPS in a pinch). The v2 MOTOR\_RF\_HAT (design-complete, unfabricated at submission) projects to $\sim$€40 for the combined controller/RF PCB and $\sim$€100 for the station core; the projection remains to be validated by fabrication.

The main open gap is reception. No protocol-valid frames were decoded from orbit across approximately 210 tracked passes; the root cause is documented in \S\ref{sec:packet-reception} and \S\ref{sec:link-budget} as the absence of an external LNA combined with a measured board-noise floor approximately 30~dB above the thermal+NF limit, and the mitigation is specified in \S\ref{sec:future-work} item 2.

The concrete deliverables of this thesis are: (a)~a built and field-validated 3D-printed AZ/EL rotator, (b)~custom RP2040 motor-controller firmware with closed-loop PID, IMU-based homing, and runaway safety, (c)~a dual-CC1200 Raspberry Pi HAT with a design-complete combined v2 variant consolidating motor control and RF front-end onto one 6-layer board, (d)~a $\sim$7\,900 LOC Python stack covering pass prediction, Doppler correction, per-pass CC1200 reconfiguration, and SatNOGS DB telemetry submission, (e)~a zero-SSH laptop-based field-deployment workflow, and (f)~a first-principles and measurement-grounded link-budget analysis that quantifies the sensitivity gap blocking orbital decode and names the specific front-end change (LNA) that would close it. The 2023 KU~Leuven masterproef by Hanssens~\cite{hanssens2023} provided the survey of the SatNOGS ecosystem and the torque-deficit finding that motivated the DC-motor path, and is referenced throughout this thesis as background rather than as a benchmark to compete with.
